\subsection{半稳定性判据}\label{src:9.5.2-subsection}
\begin{proposition}[半稳定性判据]\label{src:9.5.5}
对每个秩 $r>0$, 存在整数 $D$, 使得 $C$ 上秩 $r$, 次数 $d>D$ 的凝聚层 $F$ 半稳定, 当且仅当对所有非零凝聚子层 $E\subseteq F$, 有
\begin{equation}\label{src:9.5.6}
\frac{\mathrm h^0(E)}{\mathrm{rk}\,E}\leq\frac{\mathrm h^0(F)}{\mathrm{rk}\,F}.
\end{equation}
此外, 若 $F$ 半稳定, 则~\bookeqref{9.5.6} 中等号成立当且仅当 $\mu(E)=\mu(F)$.
\end{proposition}
我们将证明下列等价表述: 对整数 $r>0$ 和 $d$, 存在整数 $N$, 使得对每个 $n>N$ 及每个秩 $r$, 次数 $d$ 的凝聚层 $F$, 有
\begin{equation}\label{src:9.5.7}
F\text{ 半稳定}\iff\frac{\mathrm h^0(E(n))}{\mathrm{rk}\,E}\leq\frac{\mathrm h^0(F(n))}{\mathrm{rk}\,F}\quad\text{对所有 }0\neq E\subseteq F,
\end{equation}
且若 $F$ 半稳定, 则等号成立当且仅当 $\mu(E)=\mu(F)$.

回忆按定义, $F$ 半稳定当且仅当对所有 $E\subseteq F$, $\mu(E)\leq\mu(F)$. 对任意 $m\in\mathbb Z$, Riemann--Roch (\bookref{9.1.2}) 给出 $\chi(F(m))/\mathrm{rk}\,F=\mu(F)+m+1-g$, 故半稳定性等价于对所有 $0\neq E\subseteq F$, 有
\begin{equation}\label{src:9.5.8}
\frac{\chi(E(n))}{\mathrm{rk}\,E}\leq\frac{\chi(F(n))}{\mathrm{rk}\,F}.
\end{equation}
对固定的 $E$ 与 $F$, Serre 消失定理表明存在 $N$, 使得当 $n>N$ 时, $\chi(E(n))=\mathrm h^0(E(n))$, $\chi(F(n))=\mathrm h^0(F(n))$, 从而~\bookeqref{9.5.8} 等价于~\bookeqref{9.5.7} 中的不等式. 但为证明命题, 必须选取对所有 $E$ 和 $F$ 都适用的 $M$.

证明将用到实数 $x$ 的记号 $[x]_+=\max(x,0)$ 及下列界:
\begin{lemma}\label{src:9.5.9}
若 $G$ 是半稳定向量丛, 则
\[
\frac{\mathrm h^0(G)}{\mathrm{rk}\,G}\leq[\mu(G)+1]_+.
\]
\end{lemma}
\begin{proof}
若 $\mu(G)<0$, 则 $\mathrm h^0(G)=0$, 不等式成立. 对 $\mu(G)\geq0$, 按 $G$ 的次数归纳. 利用 $p\in C(\Bbbk)$ 对应的短正合列 $0\to G(-1)\to G\to G\otimes\kappa(p)\to0$, 得到 $\mu(G)=\mu(G(-1))+1$ 及 $\mathrm h^0(G)\leq\mathrm h^0(G(-1))+\mathrm{rk}\,G$. 由归纳假设 $\mathrm h^0(G(-1))/\mathrm{rk}\,G\leq\mu(G)$, 所以 $\mathrm h^0(G)/\mathrm{rk}\,G\leq\mu(G)+1$.
\end{proof}
\begin{proof}[命题~\bookref{9.5.5} 的证明]
证明~\bookeqref{9.5.7} 所述的等价性. 对~\bookeqref{9.5.7} 的 ($\Leftarrow$), 注意对商 $F\twoheadrightarrow Q$, 其中的不等式蕴含
\begin{equation}\label{src:9.5.10}
\frac{\mathrm h^0(F(n))}{r}\leq\frac{\mathrm h^0(G(n))}{\mathrm{rk}\,G}.
\end{equation}
设 $G$ 为 $F$ 的半稳定商, 且斜率 $\mu(G)\leq\mu(F)$. 将~\bookeqref{9.5.10} 与 Riemann--Roch (\bookref{9.1.2}) 及引理~\bookref{9.5.9} 结合, 得到
\[
\mu(F)+n+1-g\leq\frac{\mathrm h^0(F(n))}{r}\leq\frac{\mathrm h^0(G(n))}{\mathrm{rk}\,G}\leq[\mu(G)+n+1]_+.
\]
必要时增大 $N$, 使得对 $n>N$, $\mu(F)+n+1-g>0$. 于是 $\mu(F)-g\leq\mu(G)\leq\mu(F)$, 从而 $G$ 的秩与次数只有有限种可能. 令 $\mu_{\min}(F)$ 表示 $F$ 的 Harder--Narasimhan 滤过 (\bookref{9.1.21}) 的因子的最小斜率. 对其他任意半稳定商 $F\to Q$, 有 $\mu(Q)\geq\mu_{\min}(F)\geq\mu(F)-g$. 对所有 $E\subseteq F$ 满足 $\mu(F/E)\geq\mu(F)-g$ 的向量丛 $F$ 的轨迹有界 (习题~\bookref{9.2.9}(b)), 而秩 $\mathrm{rk}\,G$, 次数 $\deg G$ 的半稳定向量丛的轨迹也有界 (定理~\bookref{9.2.8}), 所以可增大 $N$, 使对所有 $n>N$ 及上述半稳定商 $F\to G$, 有 $\mathrm h^1(C,G(n))=0$. 因此~\bookeqref{9.5.10} 对所有半稳定商 $F\to Q$ 成立, 当且仅当对所有半稳定商 $F\to Q$, 不等式 $\mu(F)\leq\mu(Q)$ 成立. 这就证明了 $F$ 半稳定.

对~\bookeqref{9.5.7} 的 ($\Rightarrow$), 给定数 $\gamma$, 满足下列条件的对 $(E,F)$ 的轨迹有界 (习题~\bookref{9.2.9}(a)): $F$ 是秩 $r$, 次数 $d$ 的半稳定向量丛, $E\subseteq F$ 为子层且 $\mu_{\min}(E)\geq\gamma$. 由 Serre 消失定理, 可增大 $N$, 使对所有 $n>N$ 及上述对, 有 $\mathrm H^1(C,E(n))=\mathrm H^1(C,F(n))=0$. 因为 $F$ 半稳定, 对 $\mu_{\min}(E)\geq\gamma$ 的对 $(E,F)$, 不等式~\bookeqref{9.5.7} 成立.

现在选取 $\gamma:=\mu(F)-gr$, 处理 $\mu_{\min}(E)<\gamma$ 的情形. 设 $0=E_0\subseteq E_1\subseteq\cdots\subseteq E_s=E$ 是 $E$ 的 Harder--Narasimhan 滤过. 因为 $F$ 半稳定, $\mu(E_i/E_{i-1})\leq\mu(E_1)\leq\mu(F)$. 对每个因子应用引理~\bookref{9.5.9}, 得
\[
\begin{aligned}
\frac{\mathrm h^0(E(n))}{\mathrm{rk}\,E}&\leq\sum_{i=1}^s\frac{\mathrm h^0((E_i/E_{i-1})(n))}{\mathrm{rk}\,E}\\
&\leq\sum_{i=1}^s\frac{\mathrm{rk}(E_i/E_{i-1})}{\mathrm{rk}\,E}[\mu(E_i/E_{i-1})+n+1]_+\\
&\leq\sum_{i=1}^{s-1}\frac{\mathrm{rk}(E_i/E_{i-1})}{\mathrm{rk}\,E}[\mu(F)+n+1]_+\\
&\qquad+\frac{\mathrm{rk}(E/E_{s-1})}{\mathrm{rk}\,E}[\mu_{\min}(E)+n+1]_+\\
&\leq\frac{\mathrm{rk}\,E_{s-1}}{\mathrm{rk}\,E}[\mu(F)+n+1]_+\\
&\qquad+\left(1-\frac{\mathrm{rk}\,E_{s-1}}{\mathrm{rk}\,E}\right)[\mu_{\min}(E)+n+1]_+.
\end{aligned}
\]
必要时增大 $N$, 确保 $N>-(\gamma+1)$. 对 $n>N$, 有 $\mu(F)+n+1\geq0$ 及 $\mu_{\min}(E)<\gamma=\mu(F)-gr$. 因为 $1-\mathrm{rk}\,E_{s-1}/\mathrm{rk}\,E\geq1/r$, 得到
\begin{equation}\label{src:9.5.11}
\frac{\mathrm h^0(E(n))}{\mathrm{rk}\,E}<\mu(F)+n+1-g=\frac{\chi(F(n))}{\mathrm{rk}\,F}=\frac{\mathrm h^0(F(n))}{\mathrm{rk}\,F}.
\end{equation}
所以当 $\mu_{\min}(E)<\gamma$ 时, 不等式~\bookeqref{9.5.7} 也成立.

最后, 假设 $F$ 半稳定. 若子层 $E\subseteq F$ 满足 $\mu(E)=\mu(F)$, 则 $E$ 也半稳定. 因为 $E$ 的秩与次数只有有限种选择, 半稳定性的有界性 (\bookref{9.2.8}) 表明可进一步增大 $N$, 使对所有这样的 $E$ 及 $n>N$, 有 $\mathrm h^1(C,E(n))=0$. 这蕴含~\bookeqref{9.5.7} 中的不等式. 反之, 若~\bookeqref{9.5.7} 中的不等式成立, 则 $\mu_{\min}(E)\geq\gamma$, 否则不等式~\bookeqref{9.5.11} 将给出矛盾. 已经选取的 $N$ 保证对所有 $n>N$ 及 $\mu_{\min}(E)\geq\gamma$ 的 $E$, 有 $\mathrm h^1(C,E(n))=0$, 故~\bookeqref{9.5.7} 中的不等式等价于 $\mu(E)\leq\mu(F)$. 另见 \cite[命题 7.1.1 和 7.1.3]{LP97} 及 \cite[定理 4.4.1]{HL10}.
\end{proof}
